\section{Discussion}
\label{sec:discussion}

The results presented in the previous section validate the effectiveness, efficiency, and
adaptability of the DeSFAM framework in minimizing the system call attack surface and detecting
anomalous behavior in containerized environments. This section discusses key insights, design
trade-offs, and the practical implications of deploying DeSFAM in real-world container security
scenarios.

\subsection{Hybrid Analysis Advantage}
\label{subsec:hybrid_advantage}

One of the core strengths of DeSFAM lies in its hybrid syscall profiling approach. By
integrating static disassembly of container binaries with dynamic runtime tracing using eBPF
and bpftrace, DeSFAM captures both expected and runtime-dependent syscall behavior. This
dual-layer visibility ensures that rare but essential syscalls are not missed, unlike purely
static systems such as Confine~\cite{ref31_confine}, which suffer from under-approximation, or
dynamic-only systems like NIMOS~\cite{ref38_nimos}, which risk over-approximation.

\subsection{Superior Reduction Without Functional Loss}
\label{subsec:reduction_advantage}

DeSFAM achieved an average syscall reduction rate of 78.60\%, which is significantly higher than
Confine (47.60\%), Prof-gen (67.80\%), and NIMOS (63.40\%). Despite this aggressive
minimization, all target containers retained full functionality across all test workloads. This
confirms that DeSFAM's access list synthesis balances security and compatibility better than
static or log-based access list generators.

\subsection{High-Performance Real-Time Enforcement}
\label{subsec:performance_advantage}

The framework's use of eBPF for syscall filtering and policy management introduces minimal
runtime overhead. CPU usage increased by an average of only 1.74\%, and memory usage by
approximately 50~MB, well below the threshold that would impact container performance. Syscall
latency remained under 300~nanoseconds per invocation, validating DeSFAM's suitability for
high-throughput workloads such as Nginx and Redis. In comparison, user-space alternatives
(e.g., auditd) typically introduce millisecond-scale delays.

\subsection{Adaptive Response to Threats}
\label{subsec:adaptive_advantage}

DeSFAM's enforcement logic is not static. It dynamically updates syscall allowlists using risk
scoring and anomaly detection outputs from the SyscallAD module. The model achieved 96\% overall
accuracy, with an attack detection F1-score of 0.92 and a low false positive rate. This allowed
DeSFAM to detect and respond to advanced CVE-based attacks (e.g., Polkit and Dirty Pipe) in
real time, removing vulnerable syscalls such as \texttt{setuid}, \texttt{splice}, and
\texttt{clone} immediately upon detection.

\subsection{Limitations and Future Work}
\label{subsec:limitations}

While DeSFAM effectively reduces syscall exposure and detects anomalous behavior, its accuracy
remains sensitive to the representativeness of training data. Future work will incorporate
workload-aware profiling, container role classification, and continuous learning to improve
generalization across diverse workloads.

Current filtering operates at the syscall level (ID and sequence-based) without argument
inspection. Table~\ref{tab:cve_filter_rules} summarizes the contextual filtering rules applied
during CVE-based syscall extraction. Future extensions will explore argument-level analysis and
memory access tracking via eBPF to further improve detection granularity and resilience against
sophisticated attacks.

\begin{table}[htbp]
  \centering
  \caption{Filtering Rules for CVE-Based Syscall Extraction.}
  \label{tab:cve_filter_rules}
  \includegraphics[width=\linewidth]{table14}
\end{table}

Scalability in orchestrated environments is also a focus. Future deployments will integrate with
Kubernetes via CRI hooks and namespace-aware policy enforcement, ensuring consistent,
low-overhead protection across dynamic containerized clusters.
